\begin{mistake}{2026-04-12}{32}

\begin{problemcontent}
设隐函数组 \( u=u(x,y,z), v=v(x,y,z) \) 由方程组
\[ \begin{cases} x=u+vz \\ y=-u^2+v+z \end{cases} \]
确定，且 \( u(2,1,1)>0 \)。求 \( \left.\left(\frac{\partial u}{\partial x}+\frac{\partial v}{\partial y}+\frac{\partial u}{\partial z}\right)\right|_{(2,1,1)} \)。
\end{problemcontent}

\begin{solutioncontent}
1. **确定基准点数值**：将 \( (x,y,z)=(2,1,1) \) 代入原方程组：
\[ \begin{cases} 2 = u+v \\ 1 = -u^2+v+1 \end{cases} \implies \begin{cases} u+v=2 \\ v=u^2 \end{cases} \]
解得 \( u^2+u-2=0 \)，即 \( (u+2)(u-1)=0 \)。由 \( u>0 \) 知 \( u=1, v=1 \)。

2. **隐函数求导**：
对 \( x \) 求偏导：
\[ \begin{cases} 1 = u_x + v_x z \\ 0 = -2uu_x + v_x \end{cases} \xrightarrow{(1,1,1)} \begin{cases} u_x + v_x = 1 \\ -2u_x + v_x = 0 \end{cases} \implies u_x = \frac{1}{3} \]
对 \( y \) 求偏导：
\[ \begin{cases} 0 = u_y + v_y z \\ 1 = -2uu_y + v_y \end{cases} \xrightarrow{(1,1,1)} \begin{cases} u_y + v_y = 0 \\ -2u_y + v_y = 1 \end{cases} \implies v_y = \frac{2}{3} \]
对 \( z \) 求偏导（注意 \( vz \) 的乘积求导法则）：
\[ \begin{cases} 0 = u_z + (v + zv_z) \\ 0 = -2uu_z + v_z + 1 \end{cases} \xrightarrow{(1,1,1)} \begin{cases} u_z + v_z = -1 \\ -2u_z + v_z = -1 \end{cases} \implies u_z = 0 \]

3. **求和**：
\[ \left(\frac{\partial u}{\partial x}+\frac{\partial v}{\partial y}+\frac{\partial u}{\partial z}\right)\bigg|_{(2,1,1)} = \frac{1}{3} + \frac{2}{3} + 0 = 1 \]
\end{solutioncontent}

\mistakeknowledge{
\begin{itemize}
  \item \textbf{核心公式}：隐函数求导法则。对于方程组 \( F(x,y,u,v)=0, G(x,y,u,v)=0 \)，偏导数满足线性方程组，可用克莱姆法则或雅可比行列式求解。
  \item \textbf{易错点}：对 \( z \) 求导时，\( \frac{\partial(vz)}{\partial z} = v + z\frac{\partial v}{\partial z} \)，常漏掉 \( v \) 项。
  \item \textbf{关联知识}：雅可比矩阵的逆矩阵与偏导数的关系：\( \begin{pmatrix} u_x & u_y \\ v_x & v_y \end{pmatrix} = \left[\begin{pmatrix} F_u & F_v \\ G_u & G_v \end{pmatrix}\right]^{-1} \begin{pmatrix} -F_x & -F_y \\ -G_x & -G_y \end{pmatrix} \)。
\end{itemize}}

\end{mistake}
